\documentclass[12pt, a4paper]{article}

\usepackage[spanish]{babel}
\usepackage[utf8]{inputenc}
\usepackage[T1]{fontenc}
\usepackage[margin=2.5cm]{geometry}
\usepackage{graphicx}
\usepackage{amsmath, amssymb}
\usepackage{booktabs}
\usepackage{array}
\usepackage{hyperref}
\usepackage{caption}
\usepackage{subcaption}
\usepackage{float}
\usepackage{parskip}
\usepackage{setspace}
\usepackage{fancyhdr}
\usepackage{xcolor}

\hypersetup{colorlinks=true, linkcolor=blue, urlcolor=blue, citecolor=blue}
\setstretch{1.15}
\pagestyle{fancy}
\fancyhf{}
\rhead{\small Análisis Multivariante — Africa Soil Property Prediction}
\lhead{\small UPV}
\cfoot{\thepage}

\title{
  \vspace{-1cm}
  {\Large \textbf{Análisis, Monitorización y Diagnóstico de\\[4pt]
  Procesos Multivariantes}}\\[0.8cm]
  {\large Africa Soil Property Prediction Challenge:\\[2pt]
  Análisis Exploratorio, PCA y PLS sobre Espectros MIR de Suelo}\\[1cm]
  {\normalsize Universitat Politècnica de València}
}
\author{José Ramón Mena Pérez}
\date{Curso 2025--2026}

\begin{document}

\maketitle
\thispagestyle{fancy}
\vspace{1cm}

\tableofcontents
\newpage

%=============================================================================
\section{Introducción}
%=============================================================================

\subsection{Descripción del conjunto de datos}

El conjunto de datos utilizado en este trabajo procede del concurso \textit{Africa Soil
Property Prediction Challenge}, celebrado en la plataforma Kaggle en 2014 bajo el
auspicio del proyecto AfSIS (\textit{Africa Soil Information Service}). El objetivo del
concurso era demostrar que la espectroscopía infrarroja de transformada de Fourier (FTIR)
en el rango del infrarrojo medio (\textit{Mid-Infrared}, MIR), combinada con técnicas de
quimiometría multivariante, puede predecir propiedades físico-químicas del suelo de forma
rápida, no destructiva y a bajo coste, en contraposición a los métodos de análisis de
laboratorio convencionales.

El conjunto de entrenamiento contiene \textbf{1157 muestras} de suelo procedentes de
distintas localizaciones de África. Cada muestra pertenece a una de dos profundidades de
extracción: \textbf{Topsoil} (capa superficial, 0--20 cm) o \textbf{Subsoil} (capa
profunda, 20--50 cm). Un rasgo estructural clave del diseño experimental es que las
muestras están recogidas \textbf{en parejas}: en cada ubicación geográfica se extrae una
muestra de cada profundidad. Esto genera una dependencia entre observaciones que debe
tenerse en cuenta en la validación cruzada. Del total de muestras, 565 forman parejas
completas (1130 observaciones), 11 son singletons (sin pareja), y 4 ubicaciones contienen
más de dos muestras. El dataset presenta 581 muestras de Topsoil y 576 de Subsoil.

Cada muestra se caracteriza mediante tres bloques de información:

\begin{enumerate}
  \item \textbf{Espectro MIR} (3578 variables): valores de absorbancia a lo largo del
        rango espectral de 600 a 7498 cm$^{-1}$, con una resolución aproximada de
        2 cm$^{-1}$. Las variables se denominan \texttt{m600}, \texttt{m601}, \ldots,
        \texttt{m7497}.

  \item \textbf{Variables auxiliares} (15 variables): covariables geoespaciales obtenidas
        de productos de teledetección satelital, que caracterizan el entorno de cada
        punto de muestreo. No son medidas de laboratorio sino información contextual.

  \item \textbf{Variables respuesta} (5 variables): propiedades del suelo que se desea
        predecir: calcio intercambiable (\textit{Ca}), fósforo disponible (\textit{P}),
        pH del suelo, carbono orgánico del suelo (\textit{SOC}) y contenido en arena
        (\textit{Sand}).
\end{enumerate}

\subsection{Objetivos}

El presente trabajo aborda el análisis multivariante de este conjunto de datos desde una
perspectiva metodológica rigurosa, con los siguientes objetivos específicos:

\begin{enumerate}
  \item \textbf{Explorar la estructura del espacio espectral} mediante Análisis de
        Componentes Principales (PCA), identificando las direcciones de máxima varianza y
        evaluando en qué medida las primeras componentes capturan la variabilidad
        relevante del espectro MIR.

  \item \textbf{Establecer el número óptimo de componentes} mediante validación cruzada
        de diez pliegues agrupada por pares geográficos, calculando los estadísticos
        $Q^2$ (capacidad predictiva) y $R^2_X$ (varianza explicada).

  \item \textbf{Detectar anomalías y observaciones influyentes} mediante el estadístico
        $T^2$ de Hotelling (distancia en el espacio del modelo) y la Suma de Cuadrados
        de Residuos (SCR, distancia fuera del modelo), construyendo el mapa de influencia
        correspondiente.

  \item \textbf{Predecir las variables respuesta} mediante Regresión de Mínimos
        Cuadrados Parciales (PLS), tanto de forma independiente para cada target
        (PLS1) como de forma conjunta (PLS2), incluyendo las variables auxiliares
        y evaluando su contribución mediante el estadístico VIP
        (\textit{Variable Importance in Projection}).

  \item \textbf{Analizar la dificultad intrínseca del problema}, en especial la
        predicción del fósforo disponible, cuya señal espectral es indirecta, su
        distribución fuertemente asimétrica y cuya relación con el espectro MIR
        está mediada por la interacción con otros minerales del suelo.
\end{enumerate}

La importancia de este problema es significativa: disponer de modelos de calibración
espectral fiables permitiría sustituir análisis de laboratorio que tardan días o semanas
por medidas espectrales que se realizan en segundos. Esto podría transformar la gestión
agrícola en regiones donde los recursos analíticos son limitados.

%=============================================================================
\section{Análisis Exploratorio de Datos (EDA)}
%=============================================================================

\subsection{Variables espectrales: el espectro MIR}

La espectroscopía infrarroja media con transformada de Fourier (FTIR-MIR) mide la
\textbf{absorbancia} de la muestra en función del número de onda. La absorbancia en un
número de onda determinado se define como:
%
\begin{equation}
  A(\tilde{\nu}) = \log_{10}\!\left(\frac{I_0(\tilde{\nu})}{I(\tilde{\nu})}\right),
\end{equation}
%
donde $I_0$ es la intensidad del haz de referencia e $I$ la intensidad transmitida (o
reflejada, en la modalidad DRIFTS empleada en este dataset) tras interactuar con la
muestra. Un valor elevado de absorbancia indica que la muestra absorbe fuertemente en esa
frecuencia, lo que ocurre cuando la energía del fotón coincide con alguna frecuencia
vibracional de los enlaces moleculares presentes en el suelo.

A diferencia de un espectrómetro dispersivo, el instrumento FTIR mide
\textbf{simultáneamente todo el espectro}: el espejo móvil del interferómetro Michelson
recorre un rango de distancias a velocidad constante, generando un interferograma (señal
en el dominio del tiempo). La transformada de Fourier de este interferograma produce el
espectro completo de una sola vez, lo que hace el análisis extremadamente rápido.

El espectro MIR se divide convencionalmente en cuatro regiones con interpretación
diferenciada:

\begin{itemize}
  \item \textbf{$>$4000 cm$^{-1}$}: sobretonos de enlaces O--H y C--H.
  \item \textbf{1800--4000 cm$^{-1}$}: grupos funcionales con bandas intensas y
        asignables (grupos carbonilo C$=$O, carbonatos, materia orgánica).
  \item \textbf{900--1800 cm$^{-1}$}: región de \textit{huella dactilar}
        (\textit{fingerprint}). En esta zona se solapan decenas de vibraciones de
        deformación y torsión que involucran a la molécula entera. El patrón resultante
        es único para cada compuesto, como una huella dactilar, lo que hace imposible
        asignar bandas individuales pero convierte la región en la más discriminativa
        para suelos (silicatos, arcillas, materia orgánica).
  \item \textbf{$<$900 cm$^{-1}$}: vibraciones de la red inorgánica (minerales de arcilla,
        cuarzo, óxidos de hierro).
\end{itemize}

La figura~\ref{fig:eda_spectrum} muestra el espectro medio de las 1157 muestras junto con
la banda de $\pm 1$ desviación típica, y la figura~\ref{fig:eda_depth} compara los
espectros medios de Topsoil y Subsoil. Se observa que la variabilidad es muy elevada en
la región de huella dactilar, reflejando la diversidad mineralógica y orgánica de los
suelos africanos. Los Topsoils presentan sistemáticamente mayor absorbancia en la región
orgánica (1800--4000 cm$^{-1}$), coherente con un mayor contenido de materia orgánica
en la capa superficial.

\begin{figure}[H]
  \centering
  \includegraphics[width=0.95\textwidth]{figures/eda_spectrum.png}
  \caption{Espectro MIR medio de las 1157 muestras (línea azul) con banda de
           $\pm 1$ desviación típica (área sombreada). Las regiones espectrales se
           distinguen por color de fondo.}
  \label{fig:eda_spectrum}
\end{figure}

\begin{figure}[H]
  \centering
  \includegraphics[width=0.95\textwidth]{figures/eda_depth_spectra.png}
  \caption{Espectro MIR medio por profundidad (Topsoil vs.\ Subsoil). Las diferencias
           en la región orgánica son consistentes con el gradiente de materia orgánica
           con la profundidad.}
  \label{fig:eda_depth}
\end{figure}

\subsection{Variables auxiliares}

Las 15 variables auxiliares provienen de seis sensores o productos satelitales distintos,
y caracterizan el entorno físico de cada punto de muestreo (tabla~\ref{tab:aux}). Son
variables de contexto geoespacial, no medidas del suelo.

\begin{table}[H]
  \centering
  \caption{Variables auxiliares agrupadas por sensor/producto.}
  \label{tab:aux}
  \begin{tabular}{lll}
    \toprule
    \textbf{Bloque} & \textbf{Variables} & \textbf{Fuente} \\
    \midrule
    Topografía    & ELEV, RELI, CTI         & SRTM \\
    Clima         & TMAP, TMFI              & WorldClim \\
    Temperatura   & LSTD, LSTN              & MODIS MOD11A2 \\
    Reflectancia  & REF1, REF2, REF3, REF7  & MODIS MOD09A1 \\
    Albedo        & BSAN, BSAS, BSAV        & MODIS MOD43A3 \\
    Vegetación    & EVI                     & MODIS MOD13A2 \\
    \bottomrule
  \end{tabular}
\end{table}

Un rasgo importante de estas variables es que las muestras de la misma ubicación
geográfica comparten \textbf{exactamente los mismos valores auxiliares}: dado que el
píxel satelital tiene una resolución de cientos de metros, un Topsoil y su Subsoil
correspondiente caen en el mismo píxel y reciben idénticos valores. Esto ha permitido
identificar las 565 parejas del dataset comparando los vectores de variables auxiliares.
La figura~\ref{fig:eda_aux} muestra la distribución estandarizada de las 15 variables.
Las variables de reflectancia y albedo presentan mayor variabilidad relativa, mientras
que ELEV y EVI tienen distribuciones más simétricas.

\textbf{Uso en los análisis:} las variables auxiliares se utilizarán únicamente en la
etapa de \textbf{PLS} como predictores adicionales. En el \textbf{PCA} se analizará
exclusivamente el espacio espectral, y las variables auxiliares se proyectarán a
posteriori como variables suplementarias (mediante correlaciones loading) para
interpretar las componentes.

\begin{figure}[H]
  \centering
  \includegraphics[width=0.95\textwidth]{figures/eda_auxiliares.png}
  \caption{Distribución estandarizada de las variables auxiliares, coloreadas por bloque
           de sensor.}
  \label{fig:eda_aux}
\end{figure}

\subsection{Variables respuesta}

Las cinco variables respuesta del dataset (\textit{Ca}, \textit{P}, \textit{pH},
\textit{SOC}, \textit{Sand}) han sido sometidas por los organizadores del concurso a un
preprocesado previo: primero una \textbf{transformación monotónica} (tipo logarítmica o
Box-Cox) para reducir la asimetría, y después una \textbf{estandarización z}
(media$\approx$0, desviación típica$\approx$1), calculada conjuntamente sobre las 1884
observaciones de entrenamiento y test del concurso.

Esto implica que los valores que observamos en el dataset \textbf{no están en unidades
físicas originales} (mg/kg, \%, etc.), sino en la escala transformada. Los predictores
espectrales y auxiliares, en cambio, \textbf{no han sido preprocesados}: se reciben en
sus unidades crudas (absorbancia para el espectro, unidades propias de cada sensor para
las auxiliares), y será responsabilidad del analista aplicar el escalado adecuado antes
de cada análisis.

\begin{figure}[H]
  \centering
  \includegraphics[width=0.95\textwidth]{figures/eda_targets.png}
  \caption{Distribución de las cinco variables respuesta tras la transformación monotónica
           y estandarización. Se indican la media y la asimetría de cada distribución.
           El fósforo (P) presenta la asimetría más extrema (7.45), evidenciando la
           presencia de valores muy elevados que dificultan su modelización.}
  \label{fig:eda_targets}
\end{figure}

Destaca el caso del \textbf{fósforo} (P), con asimetría 7.45 incluso después de la
transformación. Su señal espectral en MIR es indirecta: el P no absorbe directamente en
este rango, sino que su concentración se correlaciona débilmente con la materia orgánica
y ciertos minerales. Esto hace que sea el target más difícil de predecir, como se
confirmará en el análisis PLS.

%=============================================================================
\section{Análisis de Componentes Principales (PCA)}
%=============================================================================

\subsection{Preprocesado: autoscaling}

Antes de aplicar el PCA, se aplica un \textbf{autoscaling} (estandarización centrado y
escalado) a la matriz espectral $\mathbf{X}$ de dimensiones $1157 \times 3578$:
%
\begin{equation}
  x_{ik}^* = \frac{x_{ik} - \bar{x}_k}{s_k},
\end{equation}
%
donde $\bar{x}_k$ y $s_k$ son la media y desviación típica de la variable $k$ calculadas
sobre las 1157 observaciones. Tras el autoscaling, cada variable tiene media 0 y varianza
1, de modo que todas contribuyen por igual al análisis independientemente de su magnitud
absoluta. Dado que las variables espectrales tienen unidades de absorbancia y rangos
similares, el efecto principal del autoscaling es eliminar la línea base y el offset
instrumental, poniendo el foco en la forma del espectro.

\subsection{Criterio de selección del número de componentes}

La figura~\ref{fig:pca_kaiser} muestra la varianza explicada por las cinco primeras
componentes principales. La primera componente (PC1) recoge el \textbf{76.3\%} de la
varianza total, una proporción extraordinariamente elevada que refleja la alta
correlación entre longitudes de onda adyacentes en espectros MIR: todos los espectros
tienen una forma global similar (misma línea base, mismas bandas dominantes de cuarzo y
materia orgánica), y el PC1 captura precisamente esa tendencia común. Las componentes
PC2 y PC3 explican el 9.5\% y el 7.5\% respectivamente, recogiendo variaciones en la
forma del espectro asociadas a diferencias mineralógicas y en la profundidad de muestreo.

\begin{figure}[H]
  \centering
  \includegraphics[width=0.80\textwidth]{figures/pca_kaiser.png}
  \caption{Varianza explicada (barras, eje izquierdo) y acumulada (línea, eje derecho)
           por las cinco primeras componentes principales. El autoscaling garantiza que
           eigenvalores superiores a la unidad indican componentes que explican más
           varianza que una variable individual.}
  \label{fig:pca_kaiser}
\end{figure}

\subsection{Validación cruzada: $Q^2$ y $R^2_X$}

Para validar el modelo PCA se ha realizado una \textbf{validación cruzada de diez
pliegues agrupada por pares geográficos} (GroupKFold). El agrupamiento es esencial:
dado que cada ubicación genera un Topsoil y un Subsoil con espectros muy similares,
permitir que estén en pliegues distintos inflaría artificialmente el $Q^2$. Al agrupar
por ubicación, las parejas permanecen siempre en el mismo pliegue.

En cada pliegue, el escalado se ajusta \textbf{únicamente sobre los datos de
entrenamiento} y se aplica al conjunto de test con los parámetros del train, evitando
así el data leakage.

El estadístico $Q^2$ para $A$ componentes se calcula como:
%
\begin{equation}
  Q^2_A = 1 - \frac{\text{PRESS}_A}{\text{TSS}},
  \qquad
  \text{PRESS}_A = \sum_{i \in \text{test}} \|\mathbf{x}_i - \hat{\mathbf{x}}_i^{(A)}\|^2,
  \qquad
  \text{TSS} = \sum_{i \in \text{test}} \|\mathbf{x}_i\|^2,
\end{equation}
%
donde $\hat{\mathbf{x}}_i^{(A)} = \mathbf{t}_i^{(A)} \mathbf{P}^{(A)\top}$ es la
reconstrucción de la muestra $i$ usando $A$ componentes.

\begin{figure}[H]
  \centering
  \includegraphics[width=0.95\textwidth]{figures/pca_cv.png}
  \caption{Izquierda: $R^2_X$ medio por componente (varianza de $\mathbf{X}$ capturada)
           con barras de error ($\pm 1$ std entre pliegues). Derecha: $Q^2$ medio de
           validación cruzada con banda de $\pm 1$ std. Los resultados son muy
           estables entre pliegues, confirmando la robustez del modelo.}
  \label{fig:pca_cv}
\end{figure}

Los resultados de la tabla~\ref{tab:cv} confirman que la estructura espectral es
extremadamente estable: el $Q^2$ supera 0.96 con tan solo $A=5$ componentes,
indicando que el modelo PCA con pocas componentes es suficiente para reconstruir
fielmente cualquier espectro MIR del dataset. La baja desviación típica entre pliegues
confirma la ausencia de sobreajuste.

\begin{table}[H]
  \centering
  \caption{Resultados de validación cruzada (media $\pm$ std sobre 10 pliegues).}
  \label{tab:cv}
  \begin{tabular}{ccccc}
    \toprule
    $A$ & $R^2_X$ (\%) & std & $Q^2$ & std \\
    \midrule
    1 & 76.32 & 0.219 & 0.761 & 0.028 \\
    2 & 85.80 & --    & 0.855 & 0.018 \\
    3 & 93.31 & --    & 0.932 & 0.009 \\
    4 & 94.96 & --    & 0.948 & 0.007 \\
    5 & 96.38 & --    & 0.963 & 0.005 \\
    \bottomrule
  \end{tabular}
\end{table}

\subsection{Score plot: Topsoil vs.\ Subsoil}

La figura~\ref{fig:pca_scores} muestra el diagrama de puntuaciones (score plot) en el
plano PC1--PC2 para las 1157 muestras, distinguiendo entre Topsoil (azul) y Subsoil
(rojo). Las elipses de confianza al 95\% se calculan a partir de la distribución
chi-cuadrado con 2 grados de libertad.

\begin{figure}[H]
  \centering
  \includegraphics[width=0.80\textwidth]{figures/pca_scores.png}
  \caption{Score plot PC1 vs.\ PC2. Las elipses representan la región de confianza al
           95\% para cada grupo bajo la suposición de normalidad multivariante.}
  \label{fig:pca_scores}
\end{figure}

Las elipses de Topsoil y Subsoil se solapan considerablemente a lo largo de PC1,
indicando que la primera componente no discrimina la profundidad. Sin embargo, en la
dirección PC2 existe una ligera separación: los Topsoils tienden a tener valores de PC2
más positivos, coherente con su mayor contenido en materia orgánica.

\subsection{Loading plots: interpretación espectral}

Los loadings $p_{1k}$ y $p_{2k}$ representan la contribución de cada variable espectral
a cada componente. La figura~\ref{fig:pca_loading1} muestra el loading plot clásico
$p_1$ vs.\ $p_2$ para las 3578 variables espectrales. Los puntos se colorean por región
espectral y se etiquetan cada 10 cm$^{-1}$ para facilitar la lectura. Las flechas
indican la dirección de recorrido del espectro de 7498 a 600 cm$^{-1}$.

\begin{figure}[H]
  \centering
  \includegraphics[width=0.92\textwidth]{figures/pca_loading1.png}
  \caption{Loading plot PC1 vs.\ PC2 para las variables espectrales. La trayectoria
           conecta los loadings en orden decreciente de número de onda. Las etiquetas
           aparecen cada 10 cm$^{-1}$.}
  \label{fig:pca_loading1}
\end{figure}

\subsection{Proyección de variables auxiliares}

Las variables auxiliares no forman parte del modelo PCA, pero pueden proyectarse como
\textbf{variables suplementarias} calculando su correlación loading con cada score:
%
\begin{equation}
  r(\text{aux}_j, \mathbf{t}_a) = \text{corr}(\mathbf{x}_j^{\text{aux}}, \mathbf{t}_a),
\end{equation}
%
lo que indica la correlación lineal entre la variable auxiliar y el score de la
componente $a$. Estos coeficientes se escalan proporcionalmente para representarlos en
el mismo espacio que los loadings espectrales.

\begin{figure}[H]
  \centering
  \includegraphics[width=0.92\textwidth]{figures/pca_loading2.png}
  \caption{Loading plot con variables auxiliares proyectadas (rombos coloreados por
           bloque de sensor). Las variables de temperatura superficial (LSTD, LSTN) y
           reflectancia se proyectan en cuadrantes opuestos a las de elevación y clima,
           reflejando gradientes geográficos en la distribución de los suelos.}
  \label{fig:pca_loading2}
\end{figure}

\subsection{Detección de anomalías: $T^2$ de Hotelling y SCR}

Para la detección de observaciones anómalas se utilizan dos estadísticos complementarios
con $A=5$ componentes:

\subsubsection*{$T^2$ de Hotelling}

Mide la distancia de cada observación al centroide del modelo \textbf{dentro} del
subespacio de las $A$ componentes:
%
\begin{equation}
  T^2_i = \sum_{a=1}^{A} \frac{t_{ia}^2}{\lambda_a},
\end{equation}
%
donde $\lambda_a$ es el autovalor (varianza) de la componente $a$. El límite de control
al nivel $\alpha$ se obtiene a partir de la distribución $F$:
%
\begin{equation}
  \text{UCL}_{T^2,\alpha} = \frac{A(n-1)}{n-A}\, F_{\alpha,A,n-A}.
\end{equation}

\subsubsection*{Suma de Cuadrados de Residuos (SCR)}

Mide la distancia de cada observación al subespacio del modelo, es decir, lo que el
modelo \textbf{no} captura:
%
\begin{equation}
  \text{SCR}_i = \|\mathbf{x}_i - \hat{\mathbf{x}}_i\|^2
               = \|\mathbf{x}_i - \mathbf{t}_i \mathbf{P}^\top\|^2.
\end{equation}
%
El límite de control se calcula mediante la aproximación de Jackson--Mudholkar, basada
en los autovalores residuales $\lambda_{A+1}, \ldots, \lambda_p$.

\begin{figure}[H]
  \centering
  \includegraphics[width=0.92\textwidth]{figures/pca_t2.png}
  \caption{Carta de control del estadístico $T^2$ de Hotelling ($A=5$). La línea
           naranja representa el UCL al 95\%, la roja al 99\% y la punteada el umbral
           de $3\times\text{UCL}_{99\%}$.}
  \label{fig:pca_t2}
\end{figure}

\begin{figure}[H]
  \centering
  \includegraphics[width=0.92\textwidth]{figures/pca_scr.png}
  \caption{Carta de control de la SCR (Jackson--Mudholkar, $A=5$). Se identifica un
           número reducido de observaciones con residuos anormalmente elevados,
           indicando espectros que no se ajustan bien al modelo.}
  \label{fig:pca_scr}
\end{figure}

La figura~\ref{fig:pca_outlier} combina ambos estadísticos en el \textbf{mapa de
influencia}. La interpretación de los cuadrantes es la siguiente:

\begin{itemize}
  \item $T^2$ alto, SCR bajo: la observación es extrema \textbf{dentro} del modelo
        (alto leverage), pero bien ajustada.
  \item $T^2$ bajo, SCR alto: la observación tiene estructura anómala que el modelo
        \textbf{no captura} (outlier de residuo).
  \item $T^2$ alto, SCR alto: anomalía en ambas dimensiones, la más preocupante.
\end{itemize}

\begin{figure}[H]
  \centering
  \includegraphics[width=0.75\textwidth]{figures/pca_outlier.png}
  \caption{Mapa de influencia $T^2$ vs.\ SCR. Las observaciones marcadas con estrella
           ($\bigstar$) superan $3\times\text{UCL}_{99\%}$ en al menos uno de los dos
           estadísticos. Las observaciones 293 y 294, correspondientes a la misma
           ubicación geográfica (una pareja Topsoil--Subsoil), presentan $T^2$ y SCR
           extremadamente elevados y serían candidatas a eliminación.}
  \label{fig:pca_outlier}
\end{figure}

El análisis identifica las observaciones \textbf{293 y 294} como las únicas que superan
el umbral de $3\times\text{UCL}_{99\%}$. El hecho de que sean una pareja
(misma ubicación) sugiere un problema sistemático en la recogida o el procesado de esa
muestra, más que un error aislado. Para los análisis PLS subsiguientes se evaluará el
impacto de su exclusión.

%=============================================================================
\section{Regresión PLS}
%=============================================================================

% ─── COMPLETAR EN LA SIGUIENTE ENTREGA ───────────────────────────────────────
%
%  Contenido previsto:
%   4.1  Preprocesado y validación cruzada PLS (GroupKFold por pares)
%   4.2  PLS1 independiente por target (Ca, pH, SOC, Sand) + fósforo separado
%   4.3  PLS2 conjunto (4 targets simultáneos)
%   4.4  Score plot PLS (T vs U) y relación interna
%   4.5  Loading plot y peso plot (W*)
%   4.6  VIP — identificación de variables espectrales relevantes
%   4.7  Comparativa PLS1 vs PLS2 y efecto de variables auxiliares
%
% ─────────────────────────────────────────────────────────────────────────────

\newpage

%=============================================================================
\section*{Referencias}
%=============================================================================

\begin{enumerate}
  \item Wold, S., Sjöström, M., \& Eriksson, L. (2001). PLS-regression: a basic tool
        of chemometrics. \textit{Chemometrics and Intelligent Laboratory Systems},
        58(2), 109--130.

  \item Jackson, J.E., \& Mudholkar, G.S. (1979). Control procedures for residuals
        associated with principal component analysis. \textit{Technometrics},
        21(3), 341--349.

  \item Nomikos, P., \& MacGregor, J.F. (1995). Multivariate SPC charts for monitoring
        batch processes. \textit{Technometrics}, 37(1), 41--59.

  \item Africa Soil Information Service (AfSIS). Kaggle Africa Soil Property Prediction
        Challenge, 2014.
        \url{https://www.kaggle.com/c/afsis-soil-properties}
\end{enumerate}

\end{document}
